\subsection{Neutrinos in Large Extra Dimensions and Short-Baseline $\nu_e$ Appearance --- arXiv:1708.09548}

\textbf{Authors:} Marcela Carena, Ying-Ying Li, Camila S. Machado, Pedro A. N. Machado, Carlos E. M. Wagner

\paragraph{主な主張}
bulk Dirac質量を加えたLED+では低いKK準位との混合を抑えながら高次準位の相対的重要性を高められ、最小LEDでは小さい短基線$\nu_\mu\to\nu_e$ appearanceを増大させうることを示す\cite{Carena:2017LEDplus}。

\paragraph{新規性}
bulk質量による波動関数の局在とKKスペクトルの変形を利用し、軽いニュートリノ質量と各KK準位への結合の最小模型での相関を変えた。

\paragraph{位置付け}
非最小LEDの代表的構成であり、$T^2$の形状だけでなくbulk質量も高次KKの振動パターンを変えるという比較対象になる。
